% Ablation study — isolates the contribution of \biotest's
% retrieval-augmented spec lookup (RAG, used by Phase B's
% metamorphic-relation mining) and the coverage-feedback
% iteration (the loop that re-scores seeds and steers Phase B
% between rounds). Strictly two-by-two: each component is
% disabled independently and jointly, yielding four
% configurations evaluated on the same six $(\mathrm{SUT,
% format})$ pairs and the same $32$-bug verified manifest used
% in \S\ref{sec:coverage} and \S\ref{sec:bug-bench}.

\subsection{Ablation Study}
\label{sec:ablation}

% Table III is declared at the top of the subsection (before the
% prose) and pinned top-only [!t] so this double-column float can
% claim the top of *this* page. Declared later, its earliest top
% slot is the next page, which pushed it onto the references page.
\begin{table*}[!t]
\centering
\caption{Line coverage (\%) per ablation. Mean$\pm$std across $n{=}4$ trials. Bold = full configuration.}
\label{tab:ablation-coverage}
\footnotesize
\setlength{\tabcolsep}{3.5pt}
\renewcommand{\arraystretch}{1.05}
\begin{tabular}{l c c c c c c c}
\toprule
& \multicolumn{3}{c}{\textbf{VCF}}
& & \multicolumn{3}{c}{\textbf{SAM}} \\
\cmidrule(lr){2-4} \cmidrule(lr){6-8}
Configuration & htsjdk & vcfpy & noodles & & htsjdk & biopython & seqan3 \\
\midrule
\textbf{\biotest{}} (full)
  & $\mathbf{45.7}$\,{\scriptsize$\pm1.7$}
  & $\mathbf{72.3}$\,{\scriptsize$\pm0.3$}
  & $\mathbf{37.1}$\,{\scriptsize$\pm0.9$}
  &
  & $\mathbf{30.9}$\,{\scriptsize$\pm0.1$}
  & $\mathbf{54.3}$\,{\scriptsize$\pm0.0$}
  & $\mathbf{93.7}$\,{\scriptsize$\pm0.3$} \\
\textit{w/o Phase A}
  & $23.75$\,{\scriptsize$\pm15.92$}
  & $59.02$\,{\scriptsize$\pm18.45$}
  & $25.75$\,{\scriptsize$\pm11.89$}
  &
  & $14.32$\,{\scriptsize$\pm16.57$}
  & $28.30$\,{\scriptsize$\pm16.69$}
  & $90.95$\,{\scriptsize$\pm5.97$} \\
\textit{w/o Phase D}
  & $41.50$\,{\scriptsize$\pm5.14$}
  & $72.03$\,{\scriptsize$\pm0.39$}
  & $34.35$\,{\scriptsize$\pm1.15$}
  &
  & $30.02$\,{\scriptsize$\pm1.68$}
  & $54.27$\,{\scriptsize$\pm0.05$}
  & $93.68$\,{\scriptsize$\pm0.21$} \\
\textit{w/o Phase A \& Phase D}
  & $\phantom{0}9.50$\,{\scriptsize$\pm19.00$}
  & $49.33$\,{\scriptsize$\pm20.35$}
  & $15.93$\,{\scriptsize$\pm13.36$}
  &
  & $\phantom{0}7.60$\,{\scriptsize$\pm15.20$}
  & $15.57$\,{\scriptsize$\pm25.15$}
  & $46.60$\,{\scriptsize$\pm53.81$} \\
\bottomrule
\end{tabular}
\end{table*}

We isolate the contribution of \biotest{}'s two principal
phases. \textbf{Phase A} supplies spec-derived domain
knowledge to the LLM that mines metamorphic relations --- it
combines a retrieval-augmented lookup over the parsed
specification with a structured prompt that encodes
precondition hints, transform compositions, and spec-rule
coverage targets; without Phase A the LLM sees only the bare
SUT name and an empty schema, with no grounding for which MRs
are valid for the format. \textbf{Phase D} is the
coverage-feedback loop that, between iterations of mining,
re-scores seeds against accumulated branch coverage,
synthesises new seeds along under-explored rules, and
re-steers Phase B's MR mining toward blind-spots; without
Phase D the pipeline reduces to a single mining round followed
by one execution pass, with no opportunity to recover from a
mining miss. Each phase is ablated independently and jointly,
yielding four configurations evaluated on the same six
$(\mathrm{SUT}, \mathrm{format})$ pairs and the same $32$-bug
verified manifest as
\S\ref{sec:coverage}--\S\ref{sec:bug-bench}. Mutation adequacy is omitted: the
mutant set is fixed by the external engines, so
syntactic kill rates saturate across configurations and
cross-ablation deltas fall inside the $\sim$1\,pp
within-tool noise band of \S\ref{sec:mutation-adequacy}. Following the
fuzz-testing methodology guidance of
Klees~et~al.~\citep{klees2018fuzz}, we run $n=4$ independent
trials per configuration and report
$\mathrm{mean} \pm \mathrm{std}$ across all cells. Trials are independent: for
\biotest{} (full) and the w/o-Phase-A configuration we use $4$ big-runs of $3$
cumulative reps, where each big-run's score is the rep mean.
For the feedback-disabled variants we use $4$ independent reps
with no cross-rep accumulation, since the loop that would build
on prior-rep state is the one being ablated.

% --------- Table 2: bug-bench ablation (single-column) ---------
\begin{table}[!htbp]
\centering
\caption{Real-bug detections per (SUT,~format) on the 32-bug
manifest (split as in Table~\ref{tab:bugbench-headline}). Bold
= column max.}
\label{tab:ablation-bugbench}
\scriptsize
\setlength{\tabcolsep}{3pt}
\renewcommand{\arraystretch}{1.15}
\begin{tabular*}{\columnwidth}{@{\extracolsep{\fill}}l c c c c c c@{}}
\toprule
& \multicolumn{3}{c}{\textbf{VCF}} & \multicolumn{2}{c|}{\textbf{SAM}} & \multirow{2}{*}{\textbf{Total}} \\
\cmidrule(lr){2-4} \cmidrule(lr){5-6}
Configuration & htsjdk & vcfpy & noodles & htsjdk & \multicolumn{1}{c|}{seqan3} & \\
\midrule
\textbf{\biotest{}} (full)        & $\mathbf{4}$ & $\mathbf{3}$ & $\mathbf{1}$ & $\mathbf{2}$ & \multicolumn{1}{c|}{$0$} & $\mathbf{10}$ \\
\textit{w/o Phase A}              & $4$          & $0$          & $0$          & $0$          & \multicolumn{1}{c|}{$0$} & $4$           \\
\textit{w/o Phase D}              & $4$          & $0$          & $0$          & $2$          & \multicolumn{1}{c|}{$0$} & $6$           \\
\textit{w/o Phase A \& Phase D}   & $1$          & $0$          & $0$          & $0$          & \multicolumn{1}{c|}{$0$} & $1$           \\
\bottomrule
\end{tabular*}
\end{table}

\paragraph{Findings.} The full configuration is the row-max on
every cell of both tables. The two ablations contribute
non-overlapping signal: removing the feedback loop is benign
on coverage for host-level Python/Java SUTs (deltas
$\leq 1$~pp on vcfpy, biopython, htsjdk/SAM, seqan3) but already
loses $4$ of $10$ real bugs --- the iterated seed-synthesis path
is what bridges from htsjdk over to vcfpy and noodles
regressions, and a single round of mining cannot recover
those. Removing domain knowledge is the heavier ablation: in
$1$--$3$ of $4$ trials per SUT, Phase~B mines $0$ valid MRs
because the LLM without spec context emits empty or
schema-invalid arrays under our prompt schema, producing the
bimodal coverage profile visible in
Table~\ref{tab:ablation-coverage} (mean kept low by trials anchored
at or near zero), and bug detection collapses from $10$
to $4$ because every bug outside the htsjdk/VCF
over-strict-rejection cluster requires a spec-grounded MR to
surface. Removing both is strictly worse
than either single ablation, with mean coverage below $50\%$ on
5 of 6 SUTs and only one bug detected --- confirming that
the two phases contribute complementary, non-redundant signal.
